% ============================================================
% Section 7: Related Work
% ============================================================
\section{Related Work}
\label{sec:related_work}

\widefloatbegin
\centering
\caption{Comparison of representative modeling attack methods on XOR APUFs.
CRP budgets are approximate and reported at or near the transition midpoint.}
\label{tab:related_comparison}
\resizebox{\textwidth}{\height}{%
\begin{tabular}{lcccccl}
\toprule
Method & Year & $k$-range & CRP required & Model & Notes \\
\midrule
LR (Ruhrmair~et~al.)~\cite{ruhrmair2010modeling}
  & 2010 & 1--4 & $10^4$--$10^5$ & LR & Fails for $k\geq4$ \\
CMA-ES (Becker)~\cite{becker2015cia}
  & 2015 & 1--6 & $10^5$--$10^6$ & CMA-ES & Slower than MLP \\
Split Attack (Wisiol~et~al.)~\cite{wisiol2020splitting}
  & 2020 & i (iPUF) & $\sim$$10^6$ & MLP & Breaks iPUF \\
MoPE (Fei~et~al.)~\cite{fei2024mope}
  & 2024 & 1--7 & $10^4$--$10^6$ & MLP+MoPE & Improves low-XOR \\
Chosen Challenge (Sayadi~et~al.)~\cite{sayadi2025}
  & 2025 & 1--64 & $10^2$--$10^5$ & LR & Req.\ challenge access \\
CalyPSO (Wisiol~et~al.)~\cite{calypso2024}
  & 2024 & 1--8 & $10^5$--$10^7$ & PSO & GPU-accelerated \\
This work & 2026 & 2--10 & $10^5$--$10^7$ & MLP & Probabilistic S-curve \\
\bottomrule
\end{tabular}
}
\widefloatend

\paragraph{Traditional ML attacks.}
Modeling attacks on Arbiter PUFs were first demonstrated by
Ruhrmair~et~al.~\cite{ruhrmair2010modeling}, who showed that logistic regression
with parity-encoded feature vectors achieves high accuracy on single APUFs.
This result established the linear separability of individual APUF responses and
motivated the XOR APUF as a countermeasure.  Subsequent work extended the attack
frontier through evolution strategies.  Becker~\cite{becker2015cia} applied
CMA-ES to 5-XOR and 6-XOR APUFs, achieving above-chance accuracy with
substantial CRP budgets.  The CalyPSO
framework~\cite{calypso2024} employed particle swarm optimization to attack
8-XOR APUFs, reporting $>95\%$ accuracy with approximately 8M CRPs.  These
evolution-based approaches search the weight space directly rather than learning
a parametric model.  While they established that higher XOR configurations are
attackable, each reported a single success point and none characterized the
\emph{probability} of success across the data spectrum.  Our work differs by
measuring $P(\text{success} \mid N, k)$ as a continuous function, revealing that
the transition is probabilistic rather than deterministic.

\paragraph{Deep learning attacks.}
Deep learning has become the dominant approach for XOR PUF modeling.
Santikellur and Chakraborty~\cite{santikellur2023} provided a comprehensive
survey and open-source implementation, establishing the standard 4-layer MLP
with Tanh hidden units as the de facto baseline.  Fei~et~al.~\cite{fei2024mope}
proposed the MoPE and MMoPE frameworks, which dynamically adjust model
complexity to match the PUF's XOR configuration, achieving success on up to
7-XOR with improved sample efficiency; their analysis emphasized that
architectural matching between model capacity and PUF complexity improves attack
efficiency.  Sayadi~et~al.~\cite{sayadi2025} introduced a chosen-challenge
attack strategy for XOR APUFs, demonstrating that selective challenge sampling
reduces the CRP requirement by exploiting systematic biases in the PUF response
distribution.  Wang~et~al.~\cite{wang2024c2d2} proposed the C2D2 framework
that combines contrastive learning with data augmentation to improve
generalization under limited CRP budgets.  While these works advanced the
state of the art in attack efficiency, they each frame the problem as a search
for the best possible architecture or training strategy.  Our work demonstrates
that below a data-dependent threshold, architectural innovations provide no
benefit, and above it, a standard MLP suffices.

\paragraph{Split attack and Interpose PUF.}
The split attack, introduced by Wisiol~et~al.~\cite{wisiol2020splitting}, fundamentally changed the modeling landscape by exploiting
structural rather than functional approximations.  Rather than learning the
complete $k$-XOR function, the split attack decomposes it into two smaller
subproblems, training separate models on the upper and lower halves of the XOR
chain and combining their probabilistic outputs.  This reduces the effective
XOR complexity by approximately a factor of 2, drastically lowering the CRP
requirement.  Crucially, Wisiol~et~al. also demonstrated that the Interpose PUF
(iPUF)~\cite{nguyen2019}, proposed as a split-attack-resistant alternative,
remains vulnerable: a $(3,3)$-iPUF is attackable with 3M CRPs, refuting the
claimed security equivalence to a 6-XOR APUF.  Our work extends this result in
two directions: a systematic sweep showing that all configurations up to
$(4,4)$ fall at $10^6$ or fewer CRPs with $>99\%$ accuracy, and a quantitative
effective-security characterization anchored to our XOR APUF scaling law---the
split attack reduces the equivalent XOR complexity by exactly 2, yielding a
CRP-budget gap of $\approx$5$\times$ to $\approx$21$\times$ on our fitted
scaling law---moving beyond the binary breakability statements typical of
prior work.

\paragraph{Theoretical perspectives.}
A growing body of work has examined the geometry of neural network loss
landscapes and its implications for learning.  Choromanska~et~al.~\cite{choromanska2015loss}
showed that the loss landscape of deep networks has a hierarchical structure
where low-energy regions form connected manifolds.  Dauphin~et~al.~\cite{dauphin2014identifying}
identified saddle points as the dominant optimization obstruction in
high-dimensional nonconvex problems.  The grokking
phenomenon~\cite{power2022grokking}---where validation accuracy remains near
random for many epochs before suddenly rising to perfect
generalization---shares qualitative features with our observed data-driven
phase transition, but occurs along the epoch axis rather than the data axis.
A recent systematic review~\cite{pufmodeling2026systematic} surveys over 70
PUF modeling studies and confirms that no prior work has characterized the
success probability as a function of CRP budget.  Our work bridges the gap
between these theoretical insights and empirical PUF modeling by showing that
the data-driven phase transition arises from a growing attractor basin in the
loss landscape, captured quantitatively by the S-curve framework.

In summary, the literature has focused on the binary question of whether a
given XOR configuration can be broken, typically under a single set of
experimental conditions.  We shift the question to a probabilistic one---how
does the probability of attack success scale with CRP budget?---and provide
the first systematic empirical characterization across 2-XOR through 9-XOR
configurations, revealing the architecture-independent data scaling law that
governs ML-based modeling attacks on XOR Arbiter PUFs.


